\section{The Hurwitz comparison on the same divisor}\label{sec:phase}
This section proves Theorem~\ref{thm:phase-orders} and the exact map from a spectral jet to prime periods. The derivative formula for the Hurwitz zeta function \cite[25.11.17]{dlmf} gives, locally normally for $|q|<1$,
\[
 \zeta(s,1+q)=\sum_{n\ge0}b_n(s)q^n,
 \qquad b_n(s)=\frac{(-1)^n(s)_n}{n!}\zeta(s+n).
\]
Here $(s)_n=s(s+1)\cdots(s+n-1)$. One obtains the identity first by an absolutely convergent binomial expansion for $\Re s>1$, then by continuation. Near a nontrivial zero $\rho$, it is normally convergent directly: $(s)_n/n!$ grows at most polynomially on a compact $s$-set, whereas $|q|^n$ decays geometrically and $\zeta(s+n)$ is uniformly bounded for $n\ge1$. Since $0<\Re\rho<1$, the Euler product proves $\zeta(\rho+n)\ne0$ for $n\ge1$, and no factor of $(\rho)_n$ vanishes.

The rotation average cancels exactly the coefficients whose degree is not a multiple of $d$. Thus
\[
 A_d(s,q)=\zeta(s)+b_d(s)q^d+O(q^{2d}),\qquad b_d(\rho)\ne0.
\]
Write $a=\zeta^{(m)}(\rho)/m!\ne0$. With $q^d=w^m$ and $s=\rho+wy$, division by $w^m$ produces $ay^m+b_d(\rho)+O(w)$. The leading polynomial has $m$ distinct roots. The implicit function theorem gives
\[
 s_j(q)=\rho+c e^{2\pi ij/m}q^{d/m}+O(q^{2d/m}),
 \qquad c^m=-b_d(\rho)/a.
\]
Every squared root difference has order $2d/m$ and a nonzero coefficient. Their product has order $d(m-1)$. A positive circuit in $q$ advances each label by $d$ modulo $m$, yielding $\gcd(m,d)$ cycles of length $m/\gcd(m,d)$. The polynomial coefficients and discriminant are holomorphic, because the contour Newton sums are holomorphic and Newton's identities express the coefficients in them. Between different disks, root differences have nonzero limiting values. Multiplication proves the full order $d\sum_i(m_i-1)$ and completes the proof of Theorem~\ref{thm:phase-orders}.

\subsection{The full local jet and its inverse}
For $p>1$, let $a=\log p$, fix a phase $\omega$ of modulus one, and set $u=\omega p^{\rho-1/2}$. The following explicit isomorphism compares the entire local multiplicity data:
\begin{equation}\label{eq:jet-isomorphism}
 \begin{split}
 \mathbb C[U,U^{-1}]/(U-u)^m&\longrightarrow
                 \mathbb C[X]/(X-\rho)^m,\\
 U&\longmapsto u\sum_{j=0}^{m-1}\frac{a^j(X-\rho)^j}{j!}.
 \end{split}
\end{equation}
Its inverse sends
\[
 X\longmapsto\rho+\frac1a\sum_{j=1}^{m-1}
       \frac{(-1)^{j+1}}j\left(\frac{U-u}{u}\right)^j.
\]
The two composites are the identity because the exponential and logarithm power series are inverse modulo the $m$th power of their nilpotent variable. The image of $U$ is a unit, so Laurent inversion is allowed. On the $j$th associated grade, $(U-u)^j$ maps to $(ua)^j(X-\rho)^j$. In particular, every local nilpotent order is preserved by this map.

Two distinct rational primes separate all complex centers jointly. Indeed, agreement of the two periods at $\rho$ and $\sigma$ gives
\[
 (\rho-\sigma)\log p=2\pi ik,\qquad
 (\rho-\sigma)\log q=2\pi il
\]
for integers $k,l$. Hence $p^l=q^k$, and prime valuations give $k=l=0$, then $\rho=\sigma$. Together with \eqref{eq:jet-isomorphism} and the Chinese remainder theorem, this yields the full finite algebra. Explicitly, the local kernel in $\mathbb C[U^{\pm1},V^{\pm1}]$ is
\[
 \left((U-u_i)^{m_i},\ V-v_i\sum_{j=0}^{m_i-1}
      \binom{\log q/\log p}{j}\left(\frac{U-u_i}{u_i}\right)^j\right).
\]
The kernels for distinct centers are comaximal, since their pairs of constant coordinates are distinct. Their intersection is the global kernel, and the quotient is $\prod_i\mathbb C[X]/(X-\rho_i)^{m_i}$.

For a single prime, two distinct centers can collide for at most one choice of the rational prime, by the same valuation argument. Therefore among any $\binom M2+1$ rational primes at least one separates every enclosed center. The minimum, over those primes, of the discriminant orders of the period images is exactly $d(M-r)$: local orders are preserved by the nonzero derivative $(\log p)p^{s-1/2}$, and cross-center collisions can only add nonnegative order.

\subsection{Two prime periods at finite precision}
Here the zero-location theorem makes the periods unit-modulus coordinates. On a fixed contour enclosing ordinates $\gamma_j$ with multiplicities $m_j$, define
\[
 S_{kl}=\sum_jm_je^{i(k\log p+l\log q)\gamma_j}
 =\frac1{2\pi i}\oint e^{(s-1/2)(k\log p+l\log q)}
                         \frac{\zeta'(s)}{\zeta(s)}\,ds.
\]
Let a target interval be $|\gamma-\gamma_0|\le h$, and let every other enclosed ordinate have distance in $[\Delta,H]$, where $0\le h<\Delta\le H$. Put $a=\log p$, $b=\log q$, $M=\sum_jm_j$, and
\[
 \sigma=\min_{\Delta\le x\le H}
       \max\{|e^{iax}-1|,|e^{ibx}-1|\}>0.
\]
Positivity follows from compactness and the two-prime separation just proved. For $K\ge2$ set $w_k=(K-|k|)/K^2$ for $|k|<K$. Given errors $|\widetilde S_{kl}-S_{kl}|\le\varepsilon$, form
\[
 B=\Re\sum_{|k|,|l|<K}w_kw_l
        e^{-i(ka+lb)\gamma_0}\widetilde S_{kl}.
\]
\begin{proposition}[An integer mass certificate]\label{prop:prime-rounding}
For the total multiplicity $m$ in the target interval,
\[
 |B-m|\le\varepsilon+
       \frac{M(K^2-1)(a^2+b^2)h^2}{12}
       +\frac{4M}{K^2\sigma^2}.
\]
An error bound below $1/2$ certifies $m$ by nearest-integer rounding.
\end{proposition}
\begin{proof}
The normalized Fej\'er kernel is
\[
 F_K(x)=K^{-2}\left|\sum_{j=0}^{K-1}e^{ijx}\right|^2
        =\sum_{|k|<K}w_ke^{ikx}.
\]
It lies between zero and one, and the geometric sum gives
$F_K(x)\le4/(K^2|1-e^{ix}|^2)$ away from its period. Averaging
$1-\cos((j-l)x)\le(j-l)^2x^2/2$ gives
$1-F_K(x)\le(K^2-1)x^2/12$. The exact statistic is
$\sum_jm_jF_K(a(\gamma_j-\gamma_0))F_K(b(\gamma_j-\gamma_0))$.
The deficit on the target is bounded by the sum of the two quadratic deficits. At every other point at least one phase distance is at least $\sigma$, so its product is bounded by $4/(K^2\sigma^2)$. The weights are nonnegative and sum to one in each coordinate, bounding the moment-error contribution by $\varepsilon$.
\end{proof}

This estimate can be made constructive on a known interval: its defining separation function is Lipschitz with constant $\max(a,b)$. Certified midpoint evaluations on a mesh, minus this constant times the half mesh size, give a lower bound for $\sigma$. Refinement eventually proves a positive bound for any fixed $[\Delta,H]$.

The finite models also quantify a precision requirement. A double atom at $\gamma$ and two unit atoms at $\gamma\pm\tau$ have moment difference
\[
 2|1-\cos((ka+lb)\tau)|\le(K-1)^2(a+b)^2\tau^2,
 \qquad |k|,|l|<K.
\]
The map between these two measures and their displayed finite moment arrays proves the bound directly. A fixed positive tolerance therefore admits arbitrarily close such arrays. To read the multiplicity of a single distinct ordinate from the rounded interval mass requires a certified interval containing just that distinct ordinate. The heat determinant certificate makes the same dependence visible through its cross-center powers in \eqref{eq:lead-lower}.
